\documentclass[11pt]{article}
\usepackage[margin=0.85in]{geometry}
\usepackage{amsmath,amssymb,booktabs,xurl,hyperref}
\usepackage[T1]{fontenc}
\setlength{\emergencystretch}{3em}
\title{Causes of the Two-Stage Controller's Scenario Failures}
\author{Pan Dynamics Collision Avoidance Research}
\date{October 1, 2026}
\begin{document}
\maketitle

\section{Finding and scope}
The fourteen interruptions have three distinct immediate causes. Eleven CLF
calls return numerically feasible, essentially optimal affine solutions but
are discarded because \texttt{coneprog} returns exit flag $-7$. Two PCBF
problems are inconsistent because the shifted affine reference, dynamics and
fixed local correction bounds cannot be satisfied together. The remaining
PCBF problem is blocked by hard collision rows in the completion tail.
These findings do not establish that the physical avoidance task is infeasible.

Two cases also overlap geometrically before their CLF interruption. Those
holds execute positive-PCBF-slack solutions; solving the relaxed affine
problem does not establish collision avoidance. An interruption and a physical
collision must not be counted as the same failure mechanism.

This is an offline diagnosis of the failed states in
\path{report/TWO_STAGE_SCENARIOS_20261001}, not another closed-loop campaign.
All fourteen failure flags were reproduced. The seventeen source/native files
remain identical to controller commit
\path{905c8bdb9cf4e790e9229133f091d6a05ae13538}. No production code,
configuration, admission rule or optimization architecture was changed.
The preceding campaign used exact observations, no road boundaries, a 6-mm
buffer, 50-ms holds, and 68/76 total stages at 8/15 m/s.

\section{Eleven CLF failures: numerical termination and result rejection}
The primary solve succeeds in all eleven cases. The secondary returns a
finite point with exit $-7$: a small search direction without passing the
solver's full feasibility/optimality termination test. The production code
rejects every nonpositive exit before inspecting the returned point.

Independent evaluation of every original cone, equality, linear inequality
and bound gives a maximum raw residual of $1.458\times10^{-12}$ over the
rejected points. Thus they satisfy their original affine constraints to very
high numerical accuracy. This is not a nonlinear safety check.

An independent objective lower bound also tests whether useful CLF progress
remained. Eliminate the first predicted state using the first dynamics
identity and measured initial state. The scaled next-step CLF error becomes
$G\Delta u_0+c$. Drop all subsequent constraints and retain only the first
input's original box $[\ell,h]$. Then
\[
 \underline\rho=\max\!\left(0,
 s\min_{\ell\le\Delta u_0\le h}\|G\Delta u_0+c\|_2^2
 -(1-\alpha)V_0\right)
 \le \rho^*.
\]
Here $s=\max(V_0,1)$ and $\alpha=0.01$. The two-variable box quadratic is
solved by its stationary point and edge minima. The eleven candidate CLF
slacks exceed this lower bound by at most $2.008\times10^{-7}$ in unscaled
CLF units; one difference is $-9.3\times10^{-15}$ from roundoff. Subject to
the measured residuals, the candidates are essentially globally optimal for
the saved convex problems, not merely feasible. This is floating-point
numerical evidence, not an interval-certified proof.

\begin{center}\small
\begin{tabular}{r l r r}
\toprule
Speed & Scenario & Returned CLF slack & Slack minus lower bound\\
\midrule
8 & headOn & 213.394021 & $8.94\times10^{-11}$\\
8 & acceleratingHeadOn & 25.799220 & $6.72\times10^{-11}$\\
8 & brakingLead & 2.691435 & $4.93\times10^{-14}$\\
8 & turningCrossing & 6.529016 & $2.43\times10^{-13}$\\
8 & curvedCrossing & 5.365320 & $1.58\times10^{-11}$\\
15 & headOn & 89.232021 & $1.26\times10^{-9}$\\
15 & acceleratingHeadOn & 91.894877 & $1.54\times10^{-8}$\\
15 & brakingLead & 7.037407 & $4.14\times10^{-12}$\\
15 & crossing & 4884.362715 & $1.43\times10^{-8}$\\
15 & turningCrossing & 867.063639 & $2.01\times10^{-7}$\\
15 & curvedCrossing & 2.311033 & $-9.33\times10^{-15}$\\
\bottomrule
\end{tabular}\end{center}

The full conic representation has disparate scales: the primary optimum is
preserved through a cap with a $10^{-6}$ additive tolerance, and the endpoint
core radii are $0.000244140625$ and $0.0009765625$ in the core's weighted
coordinates, not meters. For 8-m/s curvedCrossing, the rejected point has
solver primal measure $4.89\times10^{-14}$, dual measure $0.146013$, and
duality gap $6.96\times10^{-17}$. The dual test, not primal infeasibility,
prevents successful termination in this example.

Equivalent normalization of the cap row and endpoint cone changes eight of
the eleven exits to success. Switching the internal linear solver to normal
equations changes four to success. The original feasible set is unchanged
in those tests, and residuals were recomputed in original units. Successful
normalized cases have maximum original residual $2.77\times10^{-7}$.
Not every case is repaired by scaling; the 8-m/s brakingLead case remains
$-7$ across the representative rescaling probes. Relaxing the cap alone is
also not a general remedy. The evidence identifies numerical representation
and an overly exclusive success rule; it does not justify accepting arbitrary
failed solver outputs or declaring one scaling change sufficient.

Positive optimal CLF slack is allowed by the formulation. These failures are
not caused by a requirement that CLF slack equal zero. The second solve is
performed, and its objective reaches the independent lower bound.

\section{Three PCBF failures: genuinely inconsistent local constraints}
Removing both cones still leaves the linear problems infeasible. Numerical
Farkas witnesses were computed for $Az\le b$, $Ez=r$, including finite bounds:
\[
 \mu\ge0,\quad\boldsymbol{1}^{\mathsf T}\mu=1,\quad
 A^{\mathsf T}\mu+E^{\mathsf T}\nu\simeq0,\quad
 b^{\mathsf T}\mu+r^{\mathsf T}\nu<0.
\]
The three negative values are $-0.00885471$, $-0.00296672$ and $-0.774859$;
maximum witness equality residual is $4.86\times10^{-16}$. This localizes a
real incompatibility of the formulated rows, rather than relying only on the
solver's infeasibility message. It does not prove physical unavoidability.

\subsection{8 m/s crossing and curvedHeadOn: old affine states and fixed boxes}
The reference is made by shifting the previous \emph{affine} state trajectory,
not by rolling its inputs forward from the current state. Reuse only tests
finite values and basic speed/input domain conditions. The next problem has
both an initial mismatch and dynamics defects:
\[
 \Delta x_0=x_{\rm measured}-\bar x_0,\qquad
 \Delta x_{i+1}=A_i\Delta x_i+B_i\Delta u_i+
 \bigl(f_h(\bar x_i,\bar u_i)-\bar x_{i+1}\bigr).
\]
Meanwhile the fixed trust boxes require
\[
 |\Delta x_i|\le(2.5,2.5,0.25,2.5,1.5,0.75),\qquad
 |\Delta u_i|\le(0.075,0.125).
\]
The input coordinates are steering angle in radians and dimensionless braking
ratio. These are additional changes relative to the reference, not physical
actuator limits. The state coordinates have their respective physical units.

For both cases, removing every collision row, the terminal cone, or the
terminal geometry does not resolve infeasibility. Removing either the state
trust bounds or input trust bounds, while keeping physical limits, gives a
successful primary solve with near-zero safety slack. This demonstrates the
role of the local reference neighborhood.

For crossing, dynamics and trust boxes alone are already inconsistent. The
certificate uses first-five-stage input bounds and the yaw-rate bound at
state node 6; no collision or terminal row is needed. Maximum initial mismatch
is 0.145587 in the mixed state coordinates; maximum dynamics defect is
0.065833 rad/s in yaw rate at stage 4. Diagnostic removal of either the initial
mismatch or the dynamics-defect right-hand side makes the linear problem
feasible. Erasing these terms would change the model and is not a proposed fix.

For curvedHeadOn, physical limits also participate: the certificate includes
later steering trust bounds, yaw-rate bounds and the 1.0001-m/s positive-speed
floor at state node 25. Maximum dynamics defect is 0.210308 rad/s at stage 18.
Zeroing only the initial mismatch does not help; zeroing the dynamics-defect
terms does. This is a consistency defect in the reused affine tail combined
with local bounds and the model domain, not observation noise.

\subsection{15 m/s curvedHeadOn: hard tail collision geometry}
Removing terminal constraints or either/all trust boxes does not make this
problem feasible. Removing hard tail collision rows does. Its mismatch and
largest dynamics defect are only $1.37\times10^{-5}$ and $2.38\times10^{-5}$
in the mixed coordinates, so the large reference defects above are not the
mechanism in this frame.

The linear infeasibility witness includes hard collision-start rows at
stages 28 and 34, corresponding to absolute times 1.45 and 1.75 s. Their
separating normals are approximately $(-0.9910,-0.1336)$ and
$(0.9980,0.0633)$. Around the predicted encounter, the frozen reference
geometry changes from one almost longitudinal separation direction to the
opposite direction. These rows, dynamics and physical bounds form an
inconsistent local avoidance problem. PCBF slack cannot repair the conflict
because the completion tail has no safety slack.

This implicates the frozen hard collision approximation and reference
trajectory. It does not establish that choosing any lateral direction will
solve the case: two additional probes replacing all tail normals with one
consistent lane-lateral side both remain infeasible. No successful replacement
initialization or controller correction is demonstrated in this diagnosis.

\section{Collisions are a separate relaxed-execution issue}
The preceding independent plant/rectangle audit found first strict overlap at
1.098333 s in 15-m/s headOn and 1.050000 s in acceleratingHeadOn. The covering
controls have PCBF slack sums 0.0234122 and 0.00658919, respectively; their
unrelaxed affine collision margins are negative. They were accepted relaxed
solutions before the later solver interruption.

Minimizing safety violation does not require its minimum to be zero. The
implementation issues the solution after successful two-stage exits even
when the PCBF slack remains positive. Therefore those successful exits do not
mean the configured clearance was satisfied. Nonlinear prediction errors and
intersample motion can contribute as well; this diagnosis does not isolate
their share of the overlap. It makes no new continuous-time safety claim.

\section{Design implications and validation limits}
The two-stage objective ordering can be retained. The findings motivate
numerically suitable CLF representation and explicit numerical result
criteria, a reference trajectory consistent with the measured state and
prediction dynamics, and local collision constraints that do not trap the
hard tail around an unsuitable passage. None of these directions has been
implemented or validated here. Returning the recorded CLF points alone does
not establish success of any complete scenario or repair positive-slack
collisions.

Artifacts in \path{report/TWO_STAGE_FAILURE_CAUSES_20261001/} include original
residuals, analytic lower bounds, all ablations, witness supports, logs,
instrumentation scripts and hashes. Raw matrices and generated controller
copies remain in the documented local cache. All fourteen original exits,
all three negative certificates, the eleven lower-bound comparisons and
seventeen unchanged source/native hashes were checked. The technical report
was compiled and its text checked. No production unit suite or complete
scenario campaign was repeated because the task was diagnosis only.
\end{document}
